\chapter{How the Board Processes a Job}\label{ch:jobs}

This chapter tells you what the board does with a job. Use it to understand
the result of a plot, and to set the board for your work.

\section{One Source at a Time}

The board receives HP-GL on three links: USB, COM, and TCP port 9100. One
source plots at a time. A source keeps the plotter until its buffer is empty
and the job has ended:

\begin{itemize}
  \item On TCP, the job ends when the connection closes.
  \item On USB and COM, the job ends when the source is quiet for 2~s.
\end{itemize}

Then the next source that has data starts. The other sources wait: TCP and
USB stop at the full buffer, and COM gets XOFF. The board sends the replies
to output instructions (for example \cmd{OI} or \cmd{OS}) to the source that
plots.

If the last instruction of a job has no \code{;}, the board ends it when the
job ends. For example, \code{PU0,0} at the end of a file operates.

\section{The Job Optimizer}\label{sec:optimizer}

CAD drivers often send a drawing in a bad order: an outline in many short
parts, or holes after the outline. The \term{job optimizer} of the board
keeps a full job and then plots it in a better order:

\begin{enumerate}
  \item It joins strokes whose ends touch into longer paths.
  \item It draws everything inside a closed contour before the contour. Thus
        outlines come last, and a cutter does not free a part before its
        holes are cut.
  \item From the shapes that are ready, it draws the nearest next. It starts
        a closed contour at its nearest vertex, and an open path at its
        nearer end.
\end{enumerate}

On the AutoCAD parts that were tested, the optimizer decreased the pen-up
travel by 45\%--65\%.

\subsection{Which Sources the Optimizer Changes}

The setting \key{optimize} selects the sources. Add the values of the
sources that you want:

\begin{table}[htbp]
\centering
\caption{Values of \key{optimize}}\label{tab:optimize}
\begin{tabularx}{\linewidth}{@{}P{0.12\linewidth}L@{}}
\toprule
\textbf{Value} & \textbf{Source} \\
\midrule
1 & USB \\
2 & COM-0 or COM-1 \\
4 & Network (TCP 9100) \\
\bottomrule
\end{tabularx}
\end{table}

The default is 3 (USB and COM), because CAD drivers use these links. The
desktop app orders its own jobs, so the board does not change network jobs.

\subsection{What the Optimizer Changes}

\begin{itemize}
  \item The plot starts only when the job ends. While the board keeps the
        job, \cmd{status} shows \code{state=buffering}.
  \item A pen change or a speed change divides a job into groups. The groups
        keep their sequence.
  \item A page break (\cmd{PG}) ends a page. The board plots each page when
        it ends.
  \item A page can have 8000 points and 500 strokes. A larger job plots in
        the sequence that it arrives, and the remainder streams through.
  \item Arcs in a job that the optimizer changes are drawn as short lines
        (chords).
\end{itemize}

The \cmd{diag} command shows the last job: the number of strokes and paths,
and the pen-up travel before and after the optimizer.

\section{Arcs and Circles}

The HP-GL instructions for arcs and circles are \cmd{AA}, \cmd{AR},
\cmd{CI}, \cmd{EW}, and \cmd{WG}. With the pen down, the board draws each one
as a \term{true arc}: one smooth movement on the circle. The speed on an arc
is limited, so that the acceleration toward the center stays in
\key{accel_mm_s2}. The \cmd{status} command counts the arcs (\code{arcs=}).

The board draws an arc as \term{chords} (short straight lines) when:

\begin{itemize}
  \item the optimizer changes the job (\key{optimize}),
  \item the job scales X and Y differently, so the circle is an ellipse,
  \item the arc goes out of the clip window or the plot area,
  \item X and Y have different steps per millimeter, or
  \item the radius is less than a few motor steps.
\end{itemize}

Chords are 5° by default. The job can set a different angle with the chord
parameter, or set a maximum deviation with \cmd{CT1}. The minimum is 0.5°.

\section{Speed}

\begin{itemize}
  \item With the pen up, the speed is \key{travel_mm_s}.
  \item With the pen down, the speed is \key{draw_mm_s}. With \key{job_speed}~1,
        the \cmd{VS} instruction of a job sets it (in cm/s), up to \key{max_mm_s}.
        \cmd{VS} without a value sets it back to \key{draw_mm_s}.
  \item On curves, the board decreases the speed. The acceleration toward the
        center of the curve stays in \key{accel_mm_s2}. This applies to true
        arcs and to curves made of short lines.
  \item At a corner, the speed of each axis changes by no more than
        \key{corner_mm_s}.
\end{itemize}

The planner looks ahead along the segments in its queue and always plans to
stop at the last one. If the computer sends data too slowly, the machine
slows down. It does not lose its position.

\section{Coordinates and Clipping}

\begin{itemize}
  \item One plotter unit is \SI{0.025}{mm} (40 units per mm).
  \item The plot area is \key{width_mm} by \key{height_mm}, from plotter 0,0.
  \item The board moves the pen up to the edge of the plot area, not past it.
  \item A pen-down line is clipped to the clip window (\cmd{IW}) and to the
        plot area. The pen goes up where a line leaves the window and down
        where it comes back.
  \item A point that is less than half a plotter unit (\SI{0.0125}{mm})
        outside the plot area is on the edge. A drawing that touches the edge
        of the plot area plots without a pen lift.
  \item \cmd{RO} rotates the coordinates by 90°, 180°, or 270°.
\end{itemize}

\section{Pen Changes}

The board raises and lowers the pen only when the axes are stopped. After a
pen change, it waits for the settle time (\key{pen_down_ms} or
\key{pen_up_ms}). Between two pen changes, the segments flow through the
planner without a stop.

\section{Abort a Job}

These actions stop a job, drop the data that the board has kept, and raise
the pen:

\begin{itemize}
  \item The console command \cmd{abort}, or \menu{Abort} in the app.
  \item The device-control sequence \code{ESC.K} from the source that plots.
\end{itemize}

After a fault, the board also drops the job. A TCP job is closed. A serial
job is discarded until its source is quiet for 2~s, so the remainder of the
job does not plot.
